%============================================================
%  background.tex  –  Chapter 2
%============================================================
\chapter{Background}
\label{ch:background}

% ---------------------------------------------------------------
% WRITING TIPS – BACKGROUND
% ---------------------------------------------------------------
% Goal: Give the reader every mathematical and chemical concept they
%       need to follow the theory and method chapters.
%       Assume a reader at MSc level but not necessarily a specialist in
%       correlated electronic structure theory.
%
% Suggested sections and content:
%
%  2.1  The Electronic Structure Problem
%       - Born-Oppenheimer approximation.
%       - The many-body Schrödinger equation: Ĥ Ψ = E Ψ.
%       - Electron correlation: why Hartree-Fock is insufficient.
%         (Use the \% correlation energy metric.)
%
%  2.2  Hartree-Fock Theory
%       - Slater determinants; antisymmetry principle.
%       - The HF equations and Roothan-Hall equations.
%       - Restricted (RHF) vs. Unrestricted (UHF) formalism.
%       - Orbital energies ε_p, Koopmans' theorem.
%       - Brief: why UHF is used here (open-shell molecules, CH₂).
%
%  2.3  Molecular Orbital Integrals
%       - One-electron integrals: overlap S, core Hamiltonian h.
%       - Two-electron repulsion integrals (ERIs): chemists' vs.
%         physicists' notation. (This directly relates to your code!)
%         \[ (pq|rs) = \int\!\int \phi_p^*(1)\,\phi_q(1)\,
%            \frac{1}{r_{12}}\,\phi_r^*(2)\,\phi_s(2)\,d1\,d2 \]
%       - Antisymmetrised integrals: ⟨pq‖rs⟩ = (pr|qs) − (ps|qr).
%
%  2.4  Møller-Plesset Perturbation Theory (MP2)
%       - Partitioning of H into \hat{f} (Fock) + V.
%       - First-order wavefunction: MP1 amplitudes.
%         \[ t_{ij}^{ab} = -\frac{\langle ab \| ij \rangle}
%            {\varepsilon_a + \varepsilon_b - \varepsilon_i - \varepsilon_j} \]
%       - Second-order correlation energy: standard spin-orbital expression.
%       - Convergence, limitations, size-consistency.
%       - Spin-orbital formulation (since your code is unrestricted).
%       
%       CRITICAL ADDITIONS:
%       - Distinguish FULL MP2 (all virtuals) vs TRUNCATED MP2 (active virtuals only).
%         This is fundamental to understanding OVOS, which implements truncated MP2.
%       - Introduce T1, T2 amplitudes before Brueckner section: T1 = singles,
%         why non-zero T1 indicates suboptimal reference orbitals.
%       - When does MP2 work? Single-reference dominance vs multi-reference breakdown.
%       - Size-consistency property: E_corr(A+B) = E_corr(A) + E_corr(B).
%       
%       NICE-TO-HAVE:
%       - MP2 limitations: bond dissociation, strong correlation regimes.
%       - Computational scaling O(N^5) and why it motivates active space reduction.
%       - Connection: which formulas are for RHF, which for UHF.
%
%  2.5  Spin-Orbital Formalism
%       - Spatial vs. spin orbitals; alpha/beta blocks.
%       - Interleaved spin-orbital ordering (your code convention).
%       - Two-electron integrals in spin-orbital basis: block structure.
%         This demystifies your spatial_to_spin_eri() implementation.
%       
%       NICE-TO-HAVE:
%       - After constructing spin-orbital integrals, explain why re-sorting by
%         orbital energy (canonicalization) is necessary for MP2.
%
%  2.6  Orbital Spaces and Index Notation
%       - Occupied (i,j), active virtual (a,b), inactive virtual (e,f).
%       - Standard index conventions used throughout the thesis.
%         (Define a table of index labels here – very reader friendly!)
%       
%       CRITICAL ADDITIONS:
%       - Fully explain active virtual partitioning: why OVOS restricts MP2 to 
%         [occupied × active_virtuals] instead of [occupied × all_virtuals].
%       - Truncated MP2 energy formula with index ranges specified:
%         E_corr,trunc = Σ_{i>j ∈ occ, a>b ∈ active_virt} t_ij^ab <ij||ab>
%       - Justification: computational cost reduction O(N^5) → O(N'^5).
%       
%       NICE-TO-HAVE:
%       - Active space philosophy: how is N'_virt chosen in practice?
%       - Connection to computational accuracy vs. cost tradeoff.
%
%  2.7  Orbital Rotations and Stationarity (NEW SECTION 2.7 – insert before Brueckner)
%       CRITICAL: This section is ESSENTIAL for understanding OVOS optimization.
%       
%       [SECTION INTRO – 2.7]
%       "Orbital optimization in post-HF methods relies on rotating orbitals
%        to minimize a target energy functional. This section introduces the
%        mathematical framework for orbital rotations and the stationarity
%        principle that defines optimal orbitals."
%       
%       [SUBSECTION COVERAGE]
%       §Orbital Rotation Mathematics:
%       - Unitary transformations: U = exp(R) where R is antisymmetric.
%       - How rotations change orbital coefficients: C_new = C_old @ U.
%       - Effect on one-electron integrals: h_new = U† h_old U.
%       - Effect on two-electron integrals: g_new = U† g_old U (4-index).
%       - Effect on Fock matrix: F_new = U† F_old U.
%       - Why unitary: preserves orthonormality of orbitals.
%       
%       §The Stationarity Principle:
%       - Energy as functional of rotation parameters: E(R).
%       - Gradient: ∇_R E(R) = orbital gradient w.r.t. rotation.
%       - Stationarity condition: ∇_R E(R) = 0 defines optimal orbitals.
%       - Hessian: ∇²_R E(R) provides curvature information.
%       - Connection: OO-MP2 and Brueckner orbitals arise from stationarity.
%       
%       §Optimization Strategies (CRITICAL for Chapter 3):
%       - Newton's method: fast local convergence using gradient + Hessian.
%       - Trust region methods: restrict step size to region where model is accurate.
%       - Levenberg-Marquardt damping: interpolate between Newton and gradient descent.
%       - Convergence criteria: gradient norm threshold and energy change threshold.
%       - Note: OVOS uses adaptive Levenberg-Marquardt with trust radius.
%       
%       [LEAD-IN TO §2.8]
%       "The stationarity principle provides the mathematical foundation for
%        defining optimal orbitals. Brueckner orbitals, discussed next, are
%        a special class arising from a particular stationarity condition."
%       
%  2.8  Quantum Computing and Orbital Space (currently 2.7 – relabel to 2.8)
%       - Qubit mapping: Jordan-Wigner or Bravyi-Kitaev.
%       - Qubit count scales as 2×N_orb → motivation for OVOS.
%       - Mention VQE/QPE as target algorithms.
%
% Tips:
%   - Use tables and diagrams (e.g., orbital energy level diagrams).
%   - Every symbol used in the Theory chapter MUST be defined here.
%   - Avoid long derivations – save those for the Theory chapter or Appendix A.
%   - Reference seminal textbooks: Szabo & Ostlund, Helgaker et al.,
%     Jensen, Cramer.
% ---------------------------------------------------------------

This chapter provides the necessary theoretical background for understanding the electronic structure problem, Hartree-Fock theory, molecular orbital integrals, Møller-Plesset perturbation theory, and the spin-orbital formalism. It also introduces the orbital space partitioning and index notation conventions that are used consistently throughout this thesis. Finally, it briefly discusses the relevance of orbital space in the context of quantum computing algorithms for chemistry. Each section builds on the previous ones to establish a solid foundation for the development of the OVOS method in Chapter 3.

